\documentclass[10pt,a4paper]{amsart}
\usepackage{amsmath,amssymb,amsthm,graphicx}
% For commutative diagrams we use Xy-pic
\usepackage[all]{xy}

% We use 2cell for 2-commutative diagrams.
\xyoption{2cell}
% Theorem environments.
%
\theoremstyle{plain}
\newtheorem{theorem}[subsection]{Theorem}
\newtheorem{proposition}[subsection]{Proposition}
\newtheorem{lemma}[subsection]{Lemma}

\theoremstyle{definition}
\newtheorem{definition}[subsection]{Definition}
\newtheorem{example}[subsection]{Example}
\newtheorem{exercise}[subsection]{Exercise}
\newtheorem{situation}[subsection]{Situation}

\theoremstyle{remark}
\newtheorem{remark}[subsection]{Remark}
\newtheorem{remarks}[subsection]{Remarks}

\numberwithin{equation}{subsection}

% Macros
%
\def\lim{\mathop{\mathrm{lim}}\nolimits}
\def\colim{\mathop{\mathrm{colim}}\nolimits}
\def\Spec{\mathop{\mathrm{Spec}}}
\def\Hom{\mathop{\mathrm{Hom}}\nolimits}
\def\Ext{\mathop{\mathrm{Ext}}\nolimits}
\def\SheafHom{\mathop{\mathcal{H}\!\mathit{om}}\nolimits}
\def\SheafExt{\mathop{\mathcal{E}\!\mathit{xt}}\nolimits}
\def\Sch{\mathit{Sch}}
\def\Mor{\mathop{\mathrm{Mor}}\nolimits}
\def\Ob{\mathop{\mathrm{Ob}}\nolimits}
\def\Sh{\mathop{\mathit{Sh}}\nolimits}
\def\NL{\mathop{N\!L}\nolimits}
\def\CH{\mathop{\mathrm{CH}}\nolimits}
\def\proetale{{pro\text{-}\acute{e}tale}}
\def\etale{{\acute{e}tale}}
\def\QCoh{\mathit{QCoh}}
\def\Ker{\mathop{\mathrm{Ker}}}
\def\Im{\mathop{\mathrm{Im}}}
\def\Coker{\mathop{\mathrm{Coker}}}
\def\Coim{\mathop{\mathrm{Coim}}}

% Boxtimes
%
\DeclareMathSymbol{\boxtimes}{\mathbin}{AMSa}{"02}

%
% Macros for moduli stacks/spaces
%
\def\QCohstack{\mathcal{QC}\!\mathit{oh}}
\def\Cohstack{\mathcal{C}\!\mathit{oh}}
\def\Spacesstack{\mathcal{S}\!\mathit{paces}}
\def\Quotfunctor{\mathrm{Quot}}
\def\Hilbfunctor{\mathrm{Hilb}}
\def\Curvesstack{\mathcal{C}\!\mathit{urves}}
\def\Polarizedstack{\mathcal{P}\!\mathit{olarized}}
\def\Complexesstack{\mathcal{C}\!\mathit{omplexes}}
% \Pic is the operator that assigns to X its picard group, usage \Pic(X)
% \Picardstack_{X/B} denotes the Picard stack of X over B
% \Picardfunctor_{X/B} denotes the Picard functor of X over B
\def\Pic{\mathop{\mathrm{Pic}}\nolimits}
\def\Picardstack{\mathcal{P}\!\mathit{ic}}
\def\Picardfunctor{\mathrm{Pic}}
\def\Deformationcategory{\mathcal{D}\!\mathit{ef}}
\newenvironment{reference}{}{}
\newenvironment{slogan}{}{}
\newenvironment{history}{}{}
\usepackage{cleveref}
\crefname{originalsection}{original section}{original sections}
\Crefname{originalsection}{Original section}{Original sections}
\makeatletter
\newlabel{situation-simplicial-site}{{09WE}{source}{}{sourcepointer-1}{}}
\newlabel{situation-simplicial-site@cref}{{[originalsection][0][]09WE}{[1][0][]source}{}{}{}}
\makeatother
\crefname{originalsection}{original section}{original sections}
\Crefname{originalsection}{Original section}{Original sections}
\makeatletter
\newlabel{remark-augmentation-site}{{0D6Z}{source}{}{sourcepointer-2}{}}
\newlabel{remark-augmentation-site@cref}{{[originalsection][0][]0D6Z}{[1][0][]source}{}{}{}}
\makeatother
\crefname{originallemma}{original lemma}{original lemmas}
\Crefname{originallemma}{Original lemma}{Original lemmas}
\makeatletter
\newlabel{lemma-augmentation-cech-complex}{{0D79}{source}{}{sourcepointer-3}{}}
\newlabel{lemma-augmentation-cech-complex@cref}{{[originallemma][0][]0D79}{[1][0][]source}{}{}{}}
\makeatother
\crefname{originallemma}{original lemma}{original lemmas}
\Crefname{originallemma}{Original lemma}{Original lemmas}
\makeatletter
\newlabel{lemma-restriction-to-components-site}{{09WG}{source}{}{sourcepointer-4}{}}
\newlabel{lemma-restriction-to-components-site@cref}{{[originallemma][0][]09WG}{[1][0][]source}{}{}{}}
\makeatother
\crefname{originallemma}{original lemma}{original lemmas}
\Crefname{originallemma}{Original lemma}{Original lemmas}
\makeatletter
\newlabel{lemma-restriction-injective-to-component-site}{{09WH}{source}{}{sourcepointer-5}{}}
\newlabel{lemma-restriction-injective-to-component-site@cref}{{[originallemma][0][]09WH}{[1][0][]source}{}{}{}}
\makeatother
\crefname{originallemma}{original lemma}{original lemmas}
\Crefname{originallemma}{Original lemma}{Original lemmas}
\makeatletter
\newlabel{homology-lemma-first-quadrant-ss}{{0132}{source}{}{sourcepointer-6}{}}
\newlabel{homology-lemma-first-quadrant-ss@cref}{{[originallemma][0][]0132}{[1][0][]source}{}{}{}}
\makeatother
\crefname{originallemma}{original lemma}{original lemmas}
\Crefname{originallemma}{Original lemma}{Original lemmas}
\makeatletter
\newlabel{homology-lemma-double-complex-gives-resolution}{{0133}{source}{}{sourcepointer-7}{}}
\newlabel{homology-lemma-double-complex-gives-resolution@cref}{{[originallemma][0][]0133}{[1][0][]source}{}{}{}}
\makeatother
\crefname{originallemma}{original lemma}{original lemmas}
\Crefname{originallemma}{Original lemma}{Original lemmas}
\makeatletter
\newlabel{sites-cohomology-lemma-cohomology-of-open}{{03F3}{source}{}{sourcepointer-8}{}}
\newlabel{sites-cohomology-lemma-cohomology-of-open@cref}{{[originallemma][0][]03F3}{[1][0][]source}{}{}{}}
\makeatother
\crefname{originallemma}{original lemma}{original lemmas}
\Crefname{originallemma}{Original lemma}{Original lemmas}
\makeatletter
\newlabel{sites-cohomology-lemma-restrict-K-injective-to-open}{{08FI}{source}{}{sourcepointer-9}{}}
\newlabel{sites-cohomology-lemma-restrict-K-injective-to-open@cref}{{[originallemma][0][]08FI}{[1][0][]source}{}{}{}}
\makeatother
\crefname{originallemma}{original lemma}{original lemmas}
\Crefname{originallemma}{Original lemma}{Original lemmas}
\makeatletter
\newlabel{lemma-localize-compare}{{0D87}{source}{}{sourcepointer-10}{}}
\newlabel{lemma-localize-compare@cref}{{[originallemma][0][]0D87}{[1][0][]source}{}{}{}}
\makeatother
\crefname{originallemma}{original lemma}{original lemmas}
\Crefname{originallemma}{Original lemma}{Original lemmas}
\makeatletter
\newlabel{sites-cohomology-lemma-cohomology-on-sheaf-sets}{{07A0}{source}{}{sourcepointer-11}{}}
\newlabel{sites-cohomology-lemma-cohomology-on-sheaf-sets@cref}{{[originallemma][0][]07A0}{[1][0][]source}{}{}{}}
\makeatother
\crefname{originallemma}{original lemma}{original lemmas}
\Crefname{originallemma}{Original lemma}{Original lemmas}
\makeatletter
\newlabel{lemma-augmentation-cartesian-module}{{0D7H}{source}{}{sourcepointer-12}{}}
\newlabel{lemma-augmentation-cartesian-module@cref}{{[originallemma][0][]0D7H}{[1][0][]source}{}{}{}}
\makeatother
\crefname{originallemma}{original lemma}{original lemmas}
\Crefname{originallemma}{Original lemma}{Original lemmas}
\makeatletter
\newlabel{lemma-push-pull-localization}{{0D86}{source}{}{sourcepointer-13}{}}
\newlabel{lemma-push-pull-localization@cref}{{[originallemma][0][]0D86}{[1][0][]source}{}{}{}}
\makeatother
\crefname{originalsection}{original section}{original sections}
\Crefname{originalsection}{Original section}{Original sections}
\makeatletter
\newlabel{sites-section-glueing-sheaves}{{04TP}{source}{}{sourcepointer-14}{}}
\newlabel{sites-section-glueing-sheaves@cref}{{[originalsection][0][]04TP}{[1][0][]source}{}{}{}}
\makeatother
\crefname{originallemma}{original lemma}{original lemmas}
\Crefname{originallemma}{Original lemma}{Original lemmas}
\makeatletter
\newlabel{lemma-characterize-cartesian}{{0D7I}{source}{}{sourcepointer-15}{}}
\newlabel{lemma-characterize-cartesian@cref}{{[originallemma][0][]0D7I}{[1][0][]source}{}{}{}}
\makeatother
\crefname{originallemma}{original lemma}{original lemmas}
\Crefname{originallemma}{Original lemma}{Original lemmas}
\makeatletter
\newlabel{sites-lemma-topos-good-site}{{03CI}{source}{}{sourcepointer-16}{}}
\newlabel{sites-lemma-topos-good-site@cref}{{[originallemma][0][]03CI}{[1][0][]source}{}{}{}}
\makeatother
\crefname{originaldefinition}{original definition}{original definitions}
\Crefname{originaldefinition}{Original definition}{Original definitions}
\makeatletter
\newlabel{simplicial-definition-face-degeneracy}{{0165}{source}{}{sourcepointer-17}{}}
\newlabel{simplicial-definition-face-degeneracy@cref}{{[originaldefinition][0][]0165}{[1][0][]source}{}{}{}}
\makeatother
\crefname{originallemma}{original lemma}{original lemmas}
\Crefname{originallemma}{Original lemma}{Original lemmas}
\makeatletter
\newlabel{simplicial-lemma-work-out}{{0187}{source}{}{sourcepointer-18}{}}
\newlabel{simplicial-lemma-work-out@cref}{{[originallemma][0][]0187}{[1][0][]source}{}{}{}}
\makeatother
\crefname{originallemma}{original lemma}{original lemmas}
\Crefname{originallemma}{Original lemma}{Original lemmas}
\makeatletter
\newlabel{sites-lemma-glue-maps}{{04TQ}{source}{}{sourcepointer-19}{}}
\newlabel{sites-lemma-glue-maps@cref}{{[originallemma][0][]04TQ}{[1][0][]source}{}{}{}}
\makeatother
\crefname{originallemma}{original lemma}{original lemmas}
\Crefname{originallemma}{Original lemma}{Original lemmas}
\makeatletter
\newlabel{sites-lemma-mapping-property-glue}{{04TS}{source}{}{sourcepointer-20}{}}
\newlabel{sites-lemma-mapping-property-glue@cref}{{[originallemma][0][]04TS}{[1][0][]source}{}{}{}}
\makeatother
\crefname{originaldefinition}{original definition}{original definitions}
\Crefname{originaldefinition}{Original definition}{Original definitions}
\makeatletter
\newlabel{hypercovering-definition-homology}{{01GB}{source}{}{sourcepointer-21}{}}
\newlabel{hypercovering-definition-homology@cref}{{[originaldefinition][0][]01GB}{[1][0][]source}{}{}{}}
\makeatother
\crefname{originallemma}{original lemma}{original lemmas}
\Crefname{originallemma}{Original lemma}{Original lemmas}
\makeatletter
\newlabel{hypercovering-lemma-acyclic-hypercover-sheaves}{{01GE}{source}{}{sourcepointer-22}{}}
\newlabel{hypercovering-lemma-acyclic-hypercover-sheaves@cref}{{[originallemma][0][]01GE}{[1][0][]source}{}{}{}}
\makeatother
\crefname{originallemma}{original lemma}{original lemmas}
\Crefname{originallemma}{Original lemma}{Original lemmas}
\makeatletter
\newlabel{derived-lemma-leray-acyclicity}{{015E}{source}{}{sourcepointer-23}{}}
\newlabel{derived-lemma-leray-acyclicity@cref}{{[originallemma][0][]015E}{[1][0][]source}{}{}{}}
\makeatother
\crefname{originallemma}{original lemma}{original lemmas}
\Crefname{originallemma}{Original lemma}{Original lemmas}
\makeatletter
\newlabel{lemma-simplicial-resolution-augmentation}{{0D78}{source}{}{sourcepointer-24}{}}
\newlabel{lemma-simplicial-resolution-augmentation@cref}{{[originallemma][0][]0D78}{[1][0][]source}{}{}{}}
\makeatother
\crefname{originalsection}{original section}{original sections}
\Crefname{originalsection}{Original section}{Original sections}
\makeatletter
\newlabel{sites-cohomology-section-internal-hom}{{08J7}{source}{}{sourcepointer-25}{}}
\newlabel{sites-cohomology-section-internal-hom@cref}{{[originalsection][0][]08J7}{[1][0][]source}{}{}{}}
\makeatother
\crefname{originallemma}{original lemma}{original lemmas}
\Crefname{originallemma}{Original lemma}{Original lemmas}
\makeatletter
\newlabel{sites-cohomology-lemma-RHom-into-K-injective}{{0A94}{source}{}{sourcepointer-26}{}}
\newlabel{sites-cohomology-lemma-RHom-into-K-injective@cref}{{[originallemma][0][]0A94}{[1][0][]source}{}{}{}}
\makeatother
\crefname{originallemma}{original lemma}{original lemmas}
\Crefname{originallemma}{Original lemma}{Original lemmas}
\makeatletter
\newlabel{sites-cohomology-lemma-section-RHom-over-U}{{08JA}{source}{}{sourcepointer-27}{}}
\newlabel{sites-cohomology-lemma-section-RHom-over-U@cref}{{[originallemma][0][]08JA}{[1][0][]source}{}{}{}}
\makeatother
\crefname{originalsection}{original section}{original sections}
\Crefname{originalsection}{Original section}{Original sections}
\makeatletter
\newlabel{sites-modules-remark-j-shriek-tensor}{{0EYY}{source}{}{sourcepointer-28}{}}
\newlabel{sites-modules-remark-j-shriek-tensor@cref}{{[originalsection][0][]0EYY}{[1][0][]source}{}{}{}}
\makeatother
\crefname{originallemma}{original lemma}{original lemmas}
\Crefname{originallemma}{Original lemma}{Original lemmas}
\makeatletter
\newlabel{more-algebra-lemma-prod-qis-gives-qis}{{0EYX}{source}{}{sourcepointer-29}{}}
\newlabel{more-algebra-lemma-prod-qis-gives-qis@cref}{{[originallemma][0][]0EYX}{[1][0][]source}{}{}{}}
\makeatother
\crefname{originalsection}{original section}{original sections}
\Crefname{originalsection}{Original section}{Original sections}
\makeatletter
\newlabel{simplicial-example-fibre-products-simplicial-object}{{016E}{source}{}{sourcepointer-30}{}}
\newlabel{simplicial-example-fibre-products-simplicial-object@cref}{{[originalsection][0][]016E}{[1][0][]source}{}{}{}}
\makeatother
\crefname{originalsection}{original section}{original sections}
\Crefname{originalsection}{Original section}{Original sections}
\makeatletter
\newlabel{simplicial-remark-augmentation}{{018J}{source}{}{sourcepointer-31}{}}
\newlabel{simplicial-remark-augmentation@cref}{{[originalsection][0][]018J}{[1][0][]source}{}{}{}}
\makeatother
\crefname{originallemma}{original lemma}{original lemmas}
\Crefname{originallemma}{Original lemma}{Original lemmas}
\makeatletter
\newlabel{simplicial-lemma-cosk-minus-one}{{018I}{source}{}{sourcepointer-32}{}}
\newlabel{simplicial-lemma-cosk-minus-one@cref}{{[originallemma][0][]018I}{[1][0][]source}{}{}{}}
\makeatother
\crefname{originallemma}{original lemma}{original lemmas}
\Crefname{originallemma}{Original lemma}{Original lemmas}
\makeatletter
\newlabel{lemma-compare-cosk0}{{01GC}{source}{}{sourcepointer-33}{}}
\newlabel{lemma-compare-cosk0@cref}{{[originallemma][0][]01GC}{[1][0][]source}{}{}{}}
\makeatother
\crefname{originallemma}{original lemma}{original lemmas}
\Crefname{originallemma}{Original lemma}{Original lemmas}
\makeatletter
\newlabel{lemma-acyclicity}{{01GD}{source}{}{sourcepointer-34}{}}
\newlabel{lemma-acyclicity@cref}{{[originallemma][0][]01GD}{[1][0][]source}{}{}{}}
\makeatother
\crefname{originaldefinition}{original definition}{original definitions}
\Crefname{originaldefinition}{Original definition}{Original definitions}
\makeatletter
\newlabel{hypercovering-definition-SR}{{01G0}{source}{}{sourcepointer-35}{}}
\newlabel{hypercovering-definition-SR@cref}{{[originaldefinition][0][]01G0}{[1][0][]source}{}{}{}}
\makeatother
\crefname{originalsection}{original section}{original sections}
\Crefname{originalsection}{Original section}{Original sections}
\makeatletter
\newlabel{section-semi-representable}{{09WK}{source}{}{sourcepointer-36}{}}
\newlabel{section-semi-representable@cref}{{[originalsection][0][]09WK}{[1][0][]source}{}{}{}}
\makeatother
\crefname{originalsection}{original section}{original sections}
\Crefname{originalsection}{Original section}{Original sections}
\makeatletter
\newlabel{section-simplicial-sites}{{09WB}{source}{}{sourcepointer-37}{}}
\newlabel{section-simplicial-sites@cref}{{[originalsection][0][]09WB}{[1][0][]source}{}{}{}}
\makeatother
\crefname{originallemma}{original lemma}{original lemmas}
\Crefname{originallemma}{Original lemma}{Original lemmas}
\makeatletter
\newlabel{lemma-has-P}{{0D85}{source}{}{sourcepointer-38}{}}
\newlabel{lemma-has-P@cref}{{[originallemma][0][]0D85}{[1][0][]source}{}{}{}}
\makeatother
\crefname{originallemma}{original lemma}{original lemmas}
\Crefname{originallemma}{Original lemma}{Original lemmas}
\makeatletter
\newlabel{lemma-simplicial-cocontinuous-site}{{09WD}{source}{}{sourcepointer-39}{}}
\newlabel{lemma-simplicial-cocontinuous-site@cref}{{[originallemma][0][]09WD}{[1][0][]source}{}{}{}}
\makeatother
\crefname{originallemma}{original lemma}{original lemmas}
\Crefname{originallemma}{Original lemma}{Original lemmas}
\makeatletter
\newlabel{lemma-augmentation-site}{{0D70}{source}{}{sourcepointer-40}{}}
\newlabel{lemma-augmentation-site@cref}{{[originallemma][0][]0D70}{[1][0][]source}{}{}{}}
\makeatother
\crefname{originallemma}{original lemma}{original lemmas}
\Crefname{originallemma}{Original lemma}{Original lemmas}
\makeatletter
\newlabel{sites-lemma-when-shriek}{{00XR}{source}{}{sourcepointer-41}{}}
\newlabel{sites-lemma-when-shriek@cref}{{[originallemma][0][]00XR}{[1][0][]source}{}{}{}}
\makeatother
\crefname{originallemma}{original lemma}{original lemmas}
\Crefname{originallemma}{Original lemma}{Original lemmas}
\makeatletter
\newlabel{sites-modules-lemma-g-shriek-adjoint}{{04BG}{source}{}{sourcepointer-42}{}}
\newlabel{sites-modules-lemma-g-shriek-adjoint@cref}{{[originallemma][0][]04BG}{[1][0][]source}{}{}{}}
\makeatother
\crefname{originallemma}{original lemma}{original lemmas}
\Crefname{originallemma}{Original lemma}{Original lemmas}
\makeatletter
\newlabel{lemma-sanity-check}{{0DBZ}{source}{}{sourcepointer-43}{}}
\newlabel{lemma-sanity-check@cref}{{[originallemma][0][]0DBZ}{[1][0][]source}{}{}{}}
\makeatother
\crefname{originallemma}{original lemma}{original lemmas}
\Crefname{originallemma}{Original lemma}{Original lemmas}
\makeatletter
\newlabel{sites-modules-lemma-lower-shriek-modules}{{0797}{source}{}{sourcepointer-44}{}}
\newlabel{sites-modules-lemma-lower-shriek-modules@cref}{{[originallemma][0][]0797}{[1][0][]source}{}{}{}}
\makeatother
\crefname{originallemma}{original lemma}{original lemmas}
\Crefname{originallemma}{Original lemma}{Original lemmas}
\makeatletter
\newlabel{lemma-exactness-g-shriek-modules}{{0D74}{source}{}{sourcepointer-45}{}}
\newlabel{lemma-exactness-g-shriek-modules@cref}{{[originallemma][0][]0D74}{[1][0][]source}{}{}{}}
\makeatother
\crefname{originaldefinition}{original definition}{original definitions}
\Crefname{originaldefinition}{Original definition}{Original definitions}
\makeatletter
\newlabel{hypercovering-definition-hypercovering-variant}{{09VU}{source}{}{sourcepointer-46}{}}
\newlabel{hypercovering-definition-hypercovering-variant@cref}{{[originaldefinition][0][]09VU}{[1][0][]source}{}{}{}}
\makeatother
\title{Cohomological descent for arbitrary derived objects along a fixed hypercover}
\author{The Stacks Project Authors}
\date{stacks-project:spaces-simplicial.tex:local-snapshot-20261011; GNU FDL 1.2 or later}
\begin{document}
\maketitle
\noindent\textit{Editorial note.} Despite the historical label containingbounded-abelian, the exact statement is forE inD(C) with no boundedness hypothesis. That full scope and author proof are retained.\par\medskip
\noindent\textit{Editorial note.} The ordinary hypercoverCech-colimit theorem is not supplied as a shortcut to the unbounded unit isomorphism.\par\medskip
\noindent\textit{Allowed standard prerequisites (editorial).}\begin{itemize}
\item Explicit prerequisite boundary. Sites with equalizers and fibre products, total sites of simplicial objects, K-injective resolutions, derived adjunction and product total complexes. The local sheaf-descent,Cech-complex and augmentation proofs and the hypercover acyclicity lemma are retained.
\end{itemize}\medskip
\ifdefined\hypertarget\hypertarget{sourcepointer-1}{}\fi
\noindent\textit{Source reference.} Original section 09WE: focused-seeds/AGP/sources/stacks-snapshot-20261011/tags.txt, source lines 12492--12492. Stable Stacks tag for an expressly allowed foundational or notational input; the counted construction remains the two concrete core obligations. The retained local core is not replaced by this pointer. See standard\_prerequisites.\par\medskip
\ifdefined\hypertarget\hypertarget{sourcepointer-2}{}\fi
\noindent\textit{Source reference.} Original section 0D6Z: focused-seeds/AGP/sources/stacks-snapshot-20261011/tags.txt, source lines 16539--16539. Stable Stacks tag for an expressly allowed foundational or notational input; the counted construction remains the two concrete core obligations. The retained local core is not replaced by this pointer. See standard\_prerequisites.\par\medskip
\ifdefined\hypertarget\hypertarget{sourcepointer-3}{}\fi
\noindent\textit{Source reference.} Original lemma 0D79: focused-seeds/AGP/sources/stacks-snapshot-20261011/tags.txt, source lines 16549--16549. Stable Stacks tag for an expressly allowed foundational or notational input; the counted construction remains the two concrete core obligations. The retained local core is not replaced by this pointer. See standard\_prerequisites.\par\medskip
\ifdefined\hypertarget\hypertarget{sourcepointer-4}{}\fi
\noindent\textit{Source reference.} Original lemma 09WG: focused-seeds/AGP/sources/stacks-snapshot-20261011/tags.txt, source lines 12494--12494. Stable Stacks tag for an expressly allowed foundational or notational input; the counted construction remains the two concrete core obligations. The retained local core is not replaced by this pointer. See standard\_prerequisites.\par\medskip
\ifdefined\hypertarget\hypertarget{sourcepointer-5}{}\fi
\noindent\textit{Source reference.} Original lemma 09WH: focused-seeds/AGP/sources/stacks-snapshot-20261011/tags.txt, source lines 12495--12495. Stable Stacks tag for an expressly allowed foundational or notational input; the counted construction remains the two concrete core obligations. The retained local core is not replaced by this pointer. See standard\_prerequisites.\par\medskip
\ifdefined\hypertarget\hypertarget{sourcepointer-6}{}\fi
\noindent\textit{Source reference.} Original lemma 0132: focused-seeds/AGP/sources/stacks-snapshot-20261011/tags.txt, source lines 1420--1420. Stable Stacks tag for an expressly allowed foundational or notational input; the counted construction remains the two concrete core obligations. The retained local core is not replaced by this pointer. See standard\_prerequisites.\par\medskip
\ifdefined\hypertarget\hypertarget{sourcepointer-7}{}\fi
\noindent\textit{Source reference.} Original lemma 0133: focused-seeds/AGP/sources/stacks-snapshot-20261011/tags.txt, source lines 1421--1421. Stable Stacks tag for an expressly allowed foundational or notational input; the counted construction remains the two concrete core obligations. The retained local core is not replaced by this pointer. See standard\_prerequisites.\par\medskip
\ifdefined\hypertarget\hypertarget{sourcepointer-8}{}\fi
\noindent\textit{Source reference.} Original lemma 03F3: focused-seeds/AGP/sources/stacks-snapshot-20261011/tags.txt, source lines 4478--4478. Stable Stacks tag for an expressly allowed foundational or notational input; the counted construction remains the two concrete core obligations. The retained local core is not replaced by this pointer. See standard\_prerequisites.\par\medskip
\ifdefined\hypertarget\hypertarget{sourcepointer-9}{}\fi
\noindent\textit{Source reference.} Original lemma 08FI: focused-seeds/AGP/sources/stacks-snapshot-20261011/tags.txt, source lines 10711--10711. Stable Stacks tag for an expressly allowed foundational or notational input; the counted construction remains the two concrete core obligations. The retained local core is not replaced by this pointer. See standard\_prerequisites.\par\medskip
\ifdefined\hypertarget\hypertarget{sourcepointer-10}{}\fi
\noindent\textit{Source reference.} Original lemma 0D87: focused-seeds/AGP/sources/stacks-snapshot-20261011/tags.txt, source lines 16582--16582. Stable Stacks tag for an expressly allowed foundational or notational input; the counted construction remains the two concrete core obligations. The retained local core is not replaced by this pointer. See standard\_prerequisites.\par\medskip
\ifdefined\hypertarget\hypertarget{sourcepointer-11}{}\fi
\noindent\textit{Source reference.} Original lemma 07A0: focused-seeds/AGP/sources/stacks-snapshot-20261011/tags.txt, source lines 9290--9290. Stable Stacks tag for an expressly allowed foundational or notational input; the counted construction remains the two concrete core obligations. The retained local core is not replaced by this pointer. See standard\_prerequisites.\par\medskip
\ifdefined\hypertarget\hypertarget{sourcepointer-12}{}\fi
\noindent\textit{Source reference.} Original lemma 0D7H: focused-seeds/AGP/sources/stacks-snapshot-20261011/tags.txt, source lines 16557--16557. Stable Stacks tag for an expressly allowed foundational or notational input; the counted construction remains the two concrete core obligations. The retained local core is not replaced by this pointer. See standard\_prerequisites.\par\medskip
\ifdefined\hypertarget\hypertarget{sourcepointer-13}{}\fi
\noindent\textit{Source reference.} Original lemma 0D86: focused-seeds/AGP/sources/stacks-snapshot-20261011/tags.txt, source lines 16581--16581. Stable Stacks tag for an expressly allowed foundational or notational input; the counted construction remains the two concrete core obligations. The retained local core is not replaced by this pointer. See standard\_prerequisites.\par\medskip
\ifdefined\hypertarget\hypertarget{sourcepointer-14}{}\fi
\noindent\textit{Source reference.} Original section 04TP: focused-seeds/AGP/sources/stacks-snapshot-20261011/tags.txt, source lines 6262--6262. Stable Stacks tag for an expressly allowed foundational or notational input; the counted construction remains the two concrete core obligations. The retained local core is not replaced by this pointer. See standard\_prerequisites.\par\medskip
\ifdefined\hypertarget\hypertarget{sourcepointer-15}{}\fi
\noindent\textit{Source reference.} Original lemma 0D7I: focused-seeds/AGP/sources/stacks-snapshot-20261011/tags.txt, source lines 16558--16558. Stable Stacks tag for an expressly allowed foundational or notational input; the counted construction remains the two concrete core obligations. The retained local core is not replaced by this pointer. See standard\_prerequisites.\par\medskip
\ifdefined\hypertarget\hypertarget{sourcepointer-16}{}\fi
\noindent\textit{Source reference.} Original lemma 03CI: focused-seeds/AGP/sources/stacks-snapshot-20261011/tags.txt, source lines 4385--4385. Stable Stacks tag for an expressly allowed foundational or notational input; the counted construction remains the two concrete core obligations. The retained local core is not replaced by this pointer. See standard\_prerequisites.\par\medskip
\ifdefined\hypertarget\hypertarget{sourcepointer-17}{}\fi
\noindent\textit{Source reference.} Original definition 0165: focused-seeds/AGP/sources/stacks-snapshot-20261011/tags.txt, source lines 1532--1532. Stable Stacks tag for an expressly allowed foundational or notational input; the counted construction remains the two concrete core obligations. The retained local core is not replaced by this pointer. See standard\_prerequisites.\par\medskip
\ifdefined\hypertarget\hypertarget{sourcepointer-18}{}\fi
\noindent\textit{Source reference.} Original lemma 0187: focused-seeds/AGP/sources/stacks-snapshot-20261011/tags.txt, source lines 1606--1606. Stable Stacks tag for an expressly allowed foundational or notational input; the counted construction remains the two concrete core obligations. The retained local core is not replaced by this pointer. See standard\_prerequisites.\par\medskip
\ifdefined\hypertarget\hypertarget{sourcepointer-19}{}\fi
\noindent\textit{Source reference.} Original lemma 04TQ: focused-seeds/AGP/sources/stacks-snapshot-20261011/tags.txt, source lines 6263--6263. Stable Stacks tag for an expressly allowed foundational or notational input; the counted construction remains the two concrete core obligations. The retained local core is not replaced by this pointer. See standard\_prerequisites.\par\medskip
\ifdefined\hypertarget\hypertarget{sourcepointer-20}{}\fi
\noindent\textit{Source reference.} Original lemma 04TS: focused-seeds/AGP/sources/stacks-snapshot-20261011/tags.txt, source lines 6265--6265. Stable Stacks tag for an expressly allowed foundational or notational input; the counted construction remains the two concrete core obligations. The retained local core is not replaced by this pointer. See standard\_prerequisites.\par\medskip
\ifdefined\hypertarget\hypertarget{sourcepointer-21}{}\fi
\noindent\textit{Source reference.} Original definition 01GB: focused-seeds/AGP/sources/stacks-snapshot-20261011/tags.txt, source lines 1900--1900. Stable Stacks tag for an expressly allowed foundational or notational input; the counted construction remains the two concrete core obligations. The retained local core is not replaced by this pointer. See standard\_prerequisites.\par\medskip
\ifdefined\hypertarget\hypertarget{sourcepointer-22}{}\fi
\noindent\textit{Source reference.} Original lemma 01GE: focused-seeds/AGP/sources/stacks-snapshot-20261011/tags.txt, source lines 1903--1903. Stable Stacks tag for an expressly allowed foundational or notational input; the counted construction remains the two concrete core obligations. The retained local core is not replaced by this pointer. See standard\_prerequisites.\par\medskip
\ifdefined\hypertarget\hypertarget{sourcepointer-23}{}\fi
\noindent\textit{Source reference.} Original lemma 015E: focused-seeds/AGP/sources/stacks-snapshot-20261011/tags.txt, source lines 1505--1505. Stable Stacks tag for an expressly allowed foundational or notational input; the counted construction remains the two concrete core obligations. The retained local core is not replaced by this pointer. See standard\_prerequisites.\par\medskip
\ifdefined\hypertarget\hypertarget{sourcepointer-24}{}\fi
\noindent\textit{Source reference.} Original lemma 0D78: focused-seeds/AGP/sources/stacks-snapshot-20261011/tags.txt, source lines 16548--16548. Stable Stacks tag for an expressly allowed foundational or notational input; the counted construction remains the two concrete core obligations. The retained local core is not replaced by this pointer. See standard\_prerequisites.\par\medskip
\ifdefined\hypertarget\hypertarget{sourcepointer-25}{}\fi
\noindent\textit{Source reference.} Original section 08J7: focused-seeds/AGP/sources/stacks-snapshot-20261011/tags.txt, source lines 10840--10840. Stable Stacks tag for an expressly allowed foundational or notational input; the counted construction remains the two concrete core obligations. The retained local core is not replaced by this pointer. See standard\_prerequisites.\par\medskip
\ifdefined\hypertarget\hypertarget{sourcepointer-26}{}\fi
\noindent\textit{Source reference.} Original lemma 0A94: focused-seeds/AGP/sources/stacks-snapshot-20261011/tags.txt, source lines 12937--12937. Stable Stacks tag for an expressly allowed foundational or notational input; the counted construction remains the two concrete core obligations. The retained local core is not replaced by this pointer. See standard\_prerequisites.\par\medskip
\ifdefined\hypertarget\hypertarget{sourcepointer-27}{}\fi
\noindent\textit{Source reference.} Original lemma 08JA: focused-seeds/AGP/sources/stacks-snapshot-20261011/tags.txt, source lines 10843--10843. Stable Stacks tag for an expressly allowed foundational or notational input; the counted construction remains the two concrete core obligations. The retained local core is not replaced by this pointer. See standard\_prerequisites.\par\medskip
\ifdefined\hypertarget\hypertarget{sourcepointer-28}{}\fi
\noindent\textit{Source reference.} Original section 0EYY: focused-seeds/AGP/sources/stacks-snapshot-20261011/tags.txt, source lines 18708--18708. Stable Stacks tag for an expressly allowed foundational or notational input; the counted construction remains the two concrete core obligations. The retained local core is not replaced by this pointer. See standard\_prerequisites.\par\medskip
\ifdefined\hypertarget\hypertarget{sourcepointer-29}{}\fi
\noindent\textit{Source reference.} Original lemma 0EYX: focused-seeds/AGP/sources/stacks-snapshot-20261011/tags.txt, source lines 18707--18707. Stable Stacks tag for an expressly allowed foundational or notational input; the counted construction remains the two concrete core obligations. The retained local core is not replaced by this pointer. See standard\_prerequisites.\par\medskip
\ifdefined\hypertarget\hypertarget{sourcepointer-30}{}\fi
\noindent\textit{Source reference.} Original section 016E: focused-seeds/AGP/sources/stacks-snapshot-20261011/tags.txt, source lines 1541--1541. Stable Stacks tag for an expressly allowed foundational or notational input; the counted construction remains the two concrete core obligations. The retained local core is not replaced by this pointer. See standard\_prerequisites.\par\medskip
\ifdefined\hypertarget\hypertarget{sourcepointer-31}{}\fi
\noindent\textit{Source reference.} Original section 018J: focused-seeds/AGP/sources/stacks-snapshot-20261011/tags.txt, source lines 1618--1618. Stable Stacks tag for an expressly allowed foundational or notational input; the counted construction remains the two concrete core obligations. The retained local core is not replaced by this pointer. See standard\_prerequisites.\par\medskip
\ifdefined\hypertarget\hypertarget{sourcepointer-32}{}\fi
\noindent\textit{Source reference.} Original lemma 018I: focused-seeds/AGP/sources/stacks-snapshot-20261011/tags.txt, source lines 1617--1617. Stable Stacks tag for an expressly allowed foundational or notational input; the counted construction remains the two concrete core obligations. The retained local core is not replaced by this pointer. See standard\_prerequisites.\par\medskip
\ifdefined\hypertarget\hypertarget{sourcepointer-33}{}\fi
\noindent\textit{Source reference.} Original lemma 01GC: focused-seeds/AGP/sources/stacks-snapshot-20261011/tags.txt, source lines 1901--1901. Stable Stacks tag for an expressly allowed foundational or notational input; the counted construction remains the two concrete core obligations. The retained local core is not replaced by this pointer. See standard\_prerequisites.\par\medskip
\ifdefined\hypertarget\hypertarget{sourcepointer-34}{}\fi
\noindent\textit{Source reference.} Original lemma 01GD: focused-seeds/AGP/sources/stacks-snapshot-20261011/tags.txt, source lines 1902--1902. Stable Stacks tag for an expressly allowed foundational or notational input; the counted construction remains the two concrete core obligations. The retained local core is not replaced by this pointer. See standard\_prerequisites.\par\medskip
\ifdefined\hypertarget\hypertarget{sourcepointer-35}{}\fi
\noindent\textit{Source reference.} Original definition 01G0: focused-seeds/AGP/sources/stacks-snapshot-20261011/tags.txt, source lines 1889--1889. Stable Stacks tag for an expressly allowed foundational or notational input; the counted construction remains the two concrete core obligations. The retained local core is not replaced by this pointer. See standard\_prerequisites.\par\medskip
\ifdefined\hypertarget\hypertarget{sourcepointer-36}{}\fi
\noindent\textit{Source reference.} Original section 09WK: focused-seeds/AGP/sources/stacks-snapshot-20261011/tags.txt, source lines 12498--12498. Stable Stacks tag for an expressly allowed foundational or notational input; the counted construction remains the two concrete core obligations. The retained local core is not replaced by this pointer. See standard\_prerequisites.\par\medskip
\ifdefined\hypertarget\hypertarget{sourcepointer-37}{}\fi
\noindent\textit{Source reference.} Original section 09WB: focused-seeds/AGP/sources/stacks-snapshot-20261011/tags.txt, source lines 12489--12489. Stable Stacks tag for an expressly allowed foundational or notational input; the counted construction remains the two concrete core obligations. The retained local core is not replaced by this pointer. See standard\_prerequisites.\par\medskip
\ifdefined\hypertarget\hypertarget{sourcepointer-38}{}\fi
\noindent\textit{Source reference.} Original lemma 0D85: focused-seeds/AGP/sources/stacks-snapshot-20261011/tags.txt, source lines 16580--16580. Stable Stacks tag for an expressly allowed foundational or notational input; the counted construction remains the two concrete core obligations. The retained local core is not replaced by this pointer. See standard\_prerequisites.\par\medskip
\ifdefined\hypertarget\hypertarget{sourcepointer-39}{}\fi
\noindent\textit{Source reference.} Original lemma 09WD: focused-seeds/AGP/sources/stacks-snapshot-20261011/tags.txt, source lines 12491--12491. Stable Stacks tag for an expressly allowed foundational or notational input; the counted construction remains the two concrete core obligations. The retained local core is not replaced by this pointer. See standard\_prerequisites.\par\medskip
\ifdefined\hypertarget\hypertarget{sourcepointer-40}{}\fi
\noindent\textit{Source reference.} Original lemma 0D70: focused-seeds/AGP/sources/stacks-snapshot-20261011/tags.txt, source lines 16540--16540. Stable Stacks tag for an expressly allowed foundational or notational input; the counted construction remains the two concrete core obligations. The retained local core is not replaced by this pointer. See standard\_prerequisites.\par\medskip
\ifdefined\hypertarget\hypertarget{sourcepointer-41}{}\fi
\noindent\textit{Source reference.} Original lemma 00XR: focused-seeds/AGP/sources/stacks-snapshot-20261011/tags.txt, source lines 1229--1229. Stable Stacks tag for an expressly allowed foundational or notational input; the counted construction remains the two concrete core obligations. The retained local core is not replaced by this pointer. See standard\_prerequisites.\par\medskip
\ifdefined\hypertarget\hypertarget{sourcepointer-42}{}\fi
\noindent\textit{Source reference.} Original lemma 04BG: focused-seeds/AGP/sources/stacks-snapshot-20261011/tags.txt, source lines 5651--5651. Stable Stacks tag for an expressly allowed foundational or notational input; the counted construction remains the two concrete core obligations. The retained local core is not replaced by this pointer. See standard\_prerequisites.\par\medskip
\ifdefined\hypertarget\hypertarget{sourcepointer-43}{}\fi
\noindent\textit{Source reference.} Original lemma 0DBZ: focused-seeds/AGP/sources/stacks-snapshot-20261011/tags.txt, source lines 16714--16714. Stable Stacks tag for an expressly allowed foundational or notational input; the counted construction remains the two concrete core obligations. The retained local core is not replaced by this pointer. See standard\_prerequisites.\par\medskip
\ifdefined\hypertarget\hypertarget{sourcepointer-44}{}\fi
\noindent\textit{Source reference.} Original lemma 0797: focused-seeds/AGP/sources/stacks-snapshot-20261011/tags.txt, source lines 9262--9262. Stable Stacks tag for an expressly allowed foundational or notational input; the counted construction remains the two concrete core obligations. The retained local core is not replaced by this pointer. See standard\_prerequisites.\par\medskip
\ifdefined\hypertarget\hypertarget{sourcepointer-45}{}\fi
\noindent\textit{Source reference.} Original lemma 0D74: focused-seeds/AGP/sources/stacks-snapshot-20261011/tags.txt, source lines 16544--16544. Stable Stacks tag for an expressly allowed foundational or notational input; the counted construction remains the two concrete core obligations. The retained local core is not replaced by this pointer. See standard\_prerequisites.\par\medskip
\ifdefined\hypertarget\hypertarget{sourcepointer-46}{}\fi
\noindent\textit{Source reference.} Original definition 09VU: focused-seeds/AGP/sources/stacks-snapshot-20261011/tags.txt, source lines 12472--12472. Stable Stacks tag for an expressly allowed foundational or notational input; the counted construction remains the two concrete core obligations. The retained local core is not replaced by this pointer. See standard\_prerequisites.\par\medskip
\begin{lemma}
\label{lemma-hypercovering-descent-bounded-abelian}
Let $\mathcal{C}$ be a site with equalizers and fibre products.
Let $K$ be a hypercovering. For
$E \in D(\mathcal{C})$ the map
$$
E \longrightarrow Ra_*a^{-1}E
$$
is an isomorphism.
\end{lemma}


\begin{remark}[Variant for over an object]
\label{remark-semi-representable-over-object}
Let $\mathcal{C}$ be a site. Let $X \in \Ob(\mathcal{C})$.
The category $\text{SR}(\mathcal{C}, X)$
of {\it semi-representable objects over $X$}
is defined by the formula
$\text{SR}(\mathcal{C}, X) = \text{SR}(\mathcal{C}/X)$.
See Hypercoverings, Definition \ref{hypercovering-definition-SR}.
Thus we may apply the above discussion to the site
$\mathcal{C}/X$. Briefly, the constructions above give
\begin{enumerate}
\item a site $\mathcal{C}/K$ for $K$ in $\text{SR}(\mathcal{C}, X)$,
\item a decomposition
$\Sh(\mathcal{C}/K) = \prod \Sh(\mathcal{C}/U_i)$ if $K = \{U_i/X\}$,
\item a localization functor $j : \mathcal{C}/K \to \mathcal{C}/X$,
\item a morphism $f : \Sh(\mathcal{C}/K) \to \Sh(\mathcal{C}/L)$
for $f : K \to L$ in $\text{SR}(\mathcal{C}, X)$.
\end{enumerate}
All results of this section hold in this situation by replacing
$\mathcal{C}$ everywhere by $\mathcal{C}/X$.
\end{remark}

\begin{remark}[Ringed variant]
\label{remark-semi-representable-ringed}
Let $\mathcal{C}$ be a site. Let $\mathcal{O}_\mathcal{C}$
be a sheaf of rings on $\mathcal{C}$. In this case, for any
semi-representable object $K$ of $\mathcal{C}$ the site
$\mathcal{C}/K$ is a ringed site with sheaf
of rings $\mathcal{O}_K = j^{-1}\mathcal{O}_\mathcal{C}$.
The constructions above give
\begin{enumerate}
\item a ringed site $(\mathcal{C}/K, \mathcal{O}_K)$
for $K$ in $\text{SR}(\mathcal{C})$,
\item a decomposition
$\textit{Mod}(\mathcal{O}_K) =
\prod \textit{Mod}(\mathcal{O}_{U_i})$ if $K = \{U_i\}$,
\item a localization morphism
$j : (\Sh(\mathcal{C}/K), \mathcal{O}_K) \to
(\Sh(\mathcal{C}), \mathcal{O}_\mathcal{C})$
of ringed topoi,
\item a morphism
$f : (\Sh(\mathcal{C}/K), \mathcal{O}_K) \to
(\Sh(\mathcal{C}/L), \mathcal{O}_L)$ of ringed topoi
for $f : K \to L$ in $\text{SR}(\mathcal{C})$.
\end{enumerate}
Many of the results above hold in this setting. For example, the
functor $j^*$ has an exact left adjoint
$$
j_! : \textit{Mod}(\mathcal{O}_K) \to \textit{Mod}(\mathcal{O}_\mathcal{C}),
$$
which in terms of the product decomposition given in (2) sends
$(\mathcal{F}_i)_{i \in I}$ to $\bigoplus j_{i, !}\mathcal{F}_i$.
Similarly, given $f : K \to L$ as above, the functor $f^*$ has
an exact left adjoint
$f_! : \textit{Mod}(\mathcal{O}_K) \to \textit{Mod}(\mathcal{O}_L)$.
Thus the functors $j^*$ and $f^*$ are exact, i.e.,
$j$ and $f$ are flat morphisms of ringed topoi (also follows
from the equalities $\mathcal{O}_K = j^{-1}\mathcal{O}_\mathcal{C}$
and $\mathcal{O}_K = f^{-1}\mathcal{O}_L$).
\end{remark}

\begin{remark}[Ringed variant over an object]
\label{remark-semi-representable-ringed-over-object}
Let $\mathcal{C}$ be a site. Let $\mathcal{O}_\mathcal{C}$
be a sheaf of rings on $\mathcal{C}$. Let $X \in \Ob(\mathcal{C})$
and denote $\mathcal{O}_X = \mathcal{O}_\mathcal{C}|_{\mathcal{C}/U}$.
Then we can combine the constructions given in
Remarks \ref{remark-semi-representable-over-object}
and \ref{remark-semi-representable-ringed} to get
\begin{enumerate}
\item a ringed site $(\mathcal{C}/K, \mathcal{O}_K)$
for $K$ in $\text{SR}(\mathcal{C}, X)$,
\item a decomposition
$\textit{Mod}(\mathcal{O}_K) =
\prod \textit{Mod}(\mathcal{O}_{U_i})$ if $K = \{U_i\}$,
\item a localization morphism
$j : (\Sh(\mathcal{C}/K), \mathcal{O}_K) \to
(\Sh(\mathcal{C}/X), \mathcal{O}_X)$
of ringed topoi,
\item a morphism
$f : (\Sh(\mathcal{C}/K), \mathcal{O}_K) \to
(\Sh(\mathcal{C}/L), \mathcal{O}_L)$ of ringed topoi
for $f : K \to L$ in $\text{SR}(\mathcal{C}, X)$.
\end{enumerate}
Of course all of the results mentioned in
Remark \ref{remark-semi-representable-ringed}
hold in this setting as well.
\end{remark}





\section{The site associate to a simplicial semi-representable object}
\label{section-simplicial-semi-representable}

\noindent
Let $\mathcal{C}$ be a site. Let $K$ be a simplicial object of
$\text{SR}(\mathcal{C})$. As usual, set $K_n = K([n])$ and denote
$K(\varphi) : K_n \to K_m$ the morphism associated to $\varphi : [m] \to [n]$.
By the construction in
Section \ref{section-semi-representable} we obtain a simplicial object
$n \mapsto \mathcal{C}/K_n$ in the category whose objects are sites and
whose morphisms are cocontinuous functors. In other words, we get
a gadget as in Case B of Section \ref{section-simplicial-sites}.
The functors satisfy property P by Lemma \ref{lemma-has-P}.
Hence we may apply
Lemma \ref{lemma-simplicial-cocontinuous-site}
to obtain a site $(\mathcal{C}/K)_{total}$.

\medskip\noindent
We can describe the site $(\mathcal{C}/K)_{total}$ explicitly as follows.
Say $K_n = \{U_{n, i}\}_{i \in I_n}$. For $\varphi : [m] \to [n]$
the morphism $K(\varphi) : K_n \to K_m$ is given by a map
$\alpha(\varphi) : I_n \to I_m$ and morphisms
$f_{\varphi, i} : U_{n, i} \to U_{m, \alpha(\varphi)(i)}$ for $i \in I_n$.
Then we have
\begin{enumerate}
\item an object of $(\mathcal{C}/K)_{total}$
corresponds to an object $(U/U_{n, i})$
of $\mathcal{C}/U_{n, i}$ for some $n$ and some $i \in I_n$,
\item a morphism between $U/U_{n, i}$ and $V/U_{m, j}$
is a pair $(\varphi, f)$ where $\varphi : [m] \to [n]$,
$j = \alpha(\varphi)(i)$, and $f : U \to V$ is a morphism of
$\mathcal{C}$ such that
$$
\vcenter{
\xymatrix{
U \ar[r]_f \ar[d] & V \ar[d] \\
U_{n, i} \ar[r]^-{f_{\varphi, i}} &
U_{m, j}
}
}
$$
is commutative, and
\item coverings of the object $U/U_{n, i}$ are constructed
by starting with a covering $\{f_j : U_j \to U\}$ in $\mathcal{C}$
and letting $\{(\text{id}, f_j) : U_j/U_{n, i} \to U/U_{n, i}\}$
be a covering in $(\mathcal{C}/K)_{total}$.
\end{enumerate}
All of our general theory developed for simplicial sites applies to
$(\mathcal{C}/K)_{total}$. Observe that the obvious forgetful functor
$$
j_{total} : (\mathcal{C}/K)_{total} \longrightarrow \mathcal{C}
$$
is continuous and cocontinuous. It turns out that the associated
morphism of topoi comes from an (obvious) augmentation.

\begin{lemma}
\label{lemma-augmentation-simplicial-semi-representable}
Let $\mathcal{C}$ be a site. Let $K$ be a simplicial object of
$\text{SR}(\mathcal{C})$. The localization functor
$j_0 : \mathcal{C}/K_0 \to \mathcal{C}$ defines an augmentation
$a_0 : \Sh(\mathcal{C}/K_0) \to \Sh(\mathcal{C})$, as in case (B) of
Remark \ref{remark-augmentation-site}.
The corresponding morphisms of topoi
$$
a_n : \Sh(\mathcal{C}/K_n) \longrightarrow \Sh(\mathcal{C}),\quad
a : \Sh((\mathcal{C}/K)_{total}) \longrightarrow \Sh(\mathcal{C})
$$
of Lemma \ref{lemma-augmentation-site}
are equal to the morphisms of topoi associated to the
continuous and cocontinuous localization functors
$j_n : \mathcal{C}/K_n \to \mathcal{C}$ and
$j_{total} : (\mathcal{C}/K)_{total} \to \mathcal{C}$.
\end{lemma}

\begin{proof}
This is immediate from working through the definitions.
See in particular the footnote in the proof of
Lemma \ref{lemma-augmentation-site}
for the relationship between $a$ and $j_{total}$.
\end{proof}

\begin{lemma}
\label{lemma-comparison}
With assumption and notation as in
Lemma \ref{lemma-augmentation-simplicial-semi-representable}
we have the following properties:
\begin{enumerate}
\item there is a functor
$a^{Sh}_! : \Sh((\mathcal{C}/K)_{total}) \to \Sh(\mathcal{C})$
left adjoint to $a^{-1} : \Sh(\mathcal{C}) \to \Sh((\mathcal{C}/K)_{total})$,
\item there is a functor
$a_! : \textit{Ab}((\mathcal{C}/K)_{total}) \to \textit{Ab}(\mathcal{C})$
left adjoint to
$a^{-1} : \textit{Ab}(\mathcal{C}) \to \textit{Ab}((\mathcal{C}/K)_{total})$,
\item the functor $a^{-1}$ associates to
$\mathcal{F}$ in $\Sh(\mathcal{C})$ the sheaf on $(\mathcal{C}/K)_{total}$
which in degree $n$ is equal to $a_n^{-1}\mathcal{F}$,
\item the functor $a_*$ associates to $\mathcal{G}$ in
$\textit{Ab}((\mathcal{C}/K)_{total})$ the equalizer of the two maps
$j_{0, *}\mathcal{G}_0 \to j_{1, *}\mathcal{G}_1$,
\end{enumerate}
\end{lemma}

\begin{proof}
Parts (3) and (4) hold for any augmentation of a
simplicial site, see Lemma \ref{lemma-augmentation-site}.
Parts (1) and (2) follow as $j_{total}$ is continuous and cocontinuous.
The functor $a^{Sh}_!$ is constructed in
Sites, Lemma \ref{sites-lemma-when-shriek}
and the functor $a_!$ is constructed in
Modules on Sites, Lemma
\ref{sites-modules-lemma-g-shriek-adjoint}.
\end{proof}

\begin{lemma}
\label{lemma-sanity-check-simplicial-semi-representable}
Let $\mathcal{C}$ be a site. Let $K$ be a simplicial object of
$\text{SR}(\mathcal{C})$. Let $U/U_{n, i}$ be an object of
$\mathcal{C}/K_n$. Let
$\mathcal{F} \in \textit{Ab}((\mathcal{C}/K)_{total})$.
Then
$$
H^p(U, \mathcal{F}) = H^p(U, \mathcal{F}_{n, i})
$$
where
\begin{enumerate}
\item on the left hand side $U$ is viewed as an object of
$\mathcal{C}_{total}$, and
\item on the right hand side $\mathcal{F}_{n, i}$ is the $i$th
component of the sheaf $\mathcal{F}_n$ on $\mathcal{C}/K_n$
in the decomposition $\Sh(\mathcal{C}/K_n) = \prod \Sh(\mathcal{C}/U_{n, i})$
of Section \ref{section-semi-representable}.
\end{enumerate}
\end{lemma}

\begin{proof}
This follows immediately from Lemma \ref{lemma-sanity-check}
and the product decompositions of Section \ref{section-semi-representable}.
\end{proof}

\begin{remark}[Variant for over an object]
\label{remark-augmentation-over-object}
Let $\mathcal{C}$ be a site. Let $X \in \Ob(\mathcal{C})$.
Recall that we have a category
$\text{SR}(\mathcal{C}, X) = \text{SR}(\mathcal{C}/X)$
of semi-representable objects over $X$,
see Remark \ref{remark-semi-representable-over-object}.
We may apply the above discussion to the site
$\mathcal{C}/X$. Briefly, the constructions above give
\begin{enumerate}
\item a site $(\mathcal{C}/K)_{total}$ for a simplicial $K$ object
of $\text{SR}(\mathcal{C}, X)$,
\item a localization functor
$j_{total} : (\mathcal{C}/K)_{total} \to \mathcal{C}/X$,
\item localization functors $j_n : \mathcal{C}/K_n \to \mathcal{C}/X$,
\item a morphism of topoi
$a : \Sh((\mathcal{C}/K)_{total}) \to \Sh(\mathcal{C}/X)$,
\item morphisms of topoi
$a_n : \Sh(\mathcal{C}/K_n) \to \Sh(\mathcal{C}/X)$,
\item a functor
$a^{Sh}_! : \Sh((\mathcal{C}/K)_{total}) \to \Sh(\mathcal{C}/X)$
left adjoint to $a^{-1}$, and
\item a functor
$a_! : \textit{Ab}((\mathcal{C}/K)_{total}) \to \textit{Ab}(\mathcal{C}/X)$
left adjoint to $a^{-1}$.
\end{enumerate}
All of the results of this section hold in this setting.
To prove this one replaces
the site $\mathcal{C}$ everywhere by $\mathcal{C}/X$.
\end{remark}

\begin{remark}[Ringed variant]
\label{remark-augmentation-ringed}
Let $\mathcal{C}$ be a site. Let $\mathcal{O}_\mathcal{C}$ be a sheaf of rings.
Given a simplicial semi-representable object $K$ of $\mathcal{C}$
we set $\mathcal{O} = a^{-1}\mathcal{O}_\mathcal{C}$, where $a$
is as in Lemmas \ref{lemma-augmentation-simplicial-semi-representable} and
\ref{lemma-comparison}.
The constructions above, keeping track of the sheaves of rings
as in Remark \ref{remark-semi-representable-ringed}, give
\begin{enumerate}
\item a ringed site $((\mathcal{C}/K)_{total}, \mathcal{O})$
for a simplicial $K$ object of $\text{SR}(\mathcal{C})$,
\item a morphism of ringed topoi
$a : (\Sh((\mathcal{C}/K)_{total}), \mathcal{O}) \to
(\Sh(\mathcal{C}), \mathcal{O}_\mathcal{C})$,
\item morphisms of ringed topoi
$a_n : (\Sh(\mathcal{C}/K_n), \mathcal{O}_n) \to
(\Sh(\mathcal{C}), \mathcal{O}_\mathcal{C})$,
\item a functor
$a_! : \textit{Mod}(\mathcal{O}) \to \textit{Mod}(\mathcal{O}_\mathcal{C})$
left adjoint to $a^*$.
\end{enumerate}
The functor $a_!$ exists (but in general is not exact)
because $a^{-1}\mathcal{O}_\mathcal{C} = \mathcal{O}$
and we can replace the use of
Modules on Sites, Lemma \ref{sites-modules-lemma-g-shriek-adjoint}
in the proof of Lemma \ref{lemma-comparison}
by Modules on Sites, Lemma \ref{sites-modules-lemma-lower-shriek-modules}.
As discussed in Remark \ref{remark-semi-representable-ringed}
there are exact functors
$a_{n!} : \textit{Mod}(\mathcal{O}_n) \to
\textit{Mod}(\mathcal{O}_\mathcal{C})$
left adjoint to $a_n^*$. Consequently, the morphisms $a$ and $a_n$ are flat.
Remark \ref{remark-semi-representable-ringed}
implies the morphism of ringed topoi
$f_\varphi : (\Sh(\mathcal{C}/K_n), \mathcal{O}_n) \to
(\Sh(\mathcal{C}/K_m), \mathcal{O}_m)$
for $\varphi : [m] \to [n]$ is flat and there exists an exact functor
$f_{\varphi !} : \textit{Mod}(\mathcal{O}_n) \to \textit{Mod}(\mathcal{O}_m)$
left adjoint to $f_\varphi^*$. This in turn implies that for
the flat morphism of ringed topoi
$g_n : (\Sh(\mathcal{C}/K_n), \mathcal{O}_n) \to
(\Sh((\mathcal{C}/K)_{total}), \mathcal{O})$
the functor $g_{n!} : \textit{Mod}(\mathcal{O}_n) \to
\textit{Mod}(\mathcal{O})$ left adjoint to $g_n^*$ is exact, see
Lemma \ref{lemma-exactness-g-shriek-modules}.
\end{remark}

\begin{remark}[Ringed variant over an object]
\label{remark-augmentation-ringed-over-object}
Let $\mathcal{C}$ be a site. Let $\mathcal{O}_\mathcal{C}$ be a sheaf of rings.
Let $X \in \Ob(\mathcal{C})$ and denote
$\mathcal{O}_X = \mathcal{O}_\mathcal{C}|_{\mathcal{C}/X}$.
Then we can combine the constructions given in
Remarks \ref{remark-augmentation-over-object} and
\ref{remark-augmentation-ringed}
to get
\begin{enumerate}
\item a ringed site $((\mathcal{C}/K)_{total}, \mathcal{O})$
for a simplicial $K$ object of $\text{SR}(\mathcal{C}, X)$,
\item a morphism of ringed topoi
$a : (\Sh((\mathcal{C}/K)_{total}), \mathcal{O}) \to
(\Sh(\mathcal{C}/X), \mathcal{O}_X)$,
\item morphisms of ringed topoi
$a_n : (\Sh(\mathcal{C}/K_n), \mathcal{O}_n) \to
(\Sh(\mathcal{C}/X), \mathcal{O}_X)$,
\item a functor
$a_! : \textit{Mod}(\mathcal{O}) \to \textit{Mod}(\mathcal{O}_X)$
left adjoint to $a^*$.
\end{enumerate}
Of course, all the results mentioned in
Remark \ref{remark-augmentation-ringed}
hold in this setting as well.
\end{remark}






\section{Cohomological descent for hypercoverings}
\label{section-cohomological-descent-hypercoverings}

\noindent
Let $\mathcal{C}$ be a site. In this section we assume $\mathcal{C}$
has equalizers and fibre products. We let $K$ be a hypercovering
as defined in Hypercoverings, Definition
\ref{hypercovering-definition-hypercovering-variant}. We will study
the augmentation
$$
a : \Sh((\mathcal{C}/K)_{total}) \longrightarrow \Sh(\mathcal{C})
$$
of Section \ref{section-simplicial-semi-representable}.

\begin{lemma}
\label{lemma-augmentation-spectral-sequence}
In Situation \ref{situation-simplicial-site} let
$a_0$ be an augmentation towards a site $\mathcal{D}$
as in Remark \ref{remark-augmentation-site}.
For any $K$ in $D^+(\mathcal{C}_{total})$ there is a spectral
sequence 
$(E_r, d_r)_{r \geq 0}$ with
$$
E_1^{p, q} = R^qa_{p, *} K_p,\quad
d_1^{p, q} : E_1^{p, q} \to E_1^{p + 1, q}
$$
converging to $R^{p + q}a_*K$. This spectral sequence is functorial in $K$.
\end{lemma}

\begin{proof}
Let $\mathcal{I}^\bullet$ be a bounded below complex of injectives
representing $K$. Consider the double complex with terms
$$
A^{p, q} = a_{p, *}\mathcal{I}^q_p
$$
where the horizontal arrows come from
Lemma \ref{lemma-augmentation-cech-complex}
and the vertical arrows from the differentials of the
complex $\mathcal{I}^\bullet$. The rows of the double complex are exact
in positive degrees and evaluate to $a_*\mathcal{I}^q$ in degree $0$.
On the other hand, since restriction to $\mathcal{C}_p$ is exact
(Lemma \ref{lemma-restriction-to-components-site})
the complex $\mathcal{I}_p^\bullet$ represents $K_p$ in
$D(\mathcal{C}_p)$. The sheaves $\mathcal{I}_p^q$ are injective
abelian sheaves on $\mathcal{C}_p$
(Lemma \ref{lemma-restriction-injective-to-component-site}).
Hence the cohomology of the columns computes $R^qa_{p, *}K_p$.
We conclude by applying
Homology, Lemmas \ref{homology-lemma-first-quadrant-ss} and
\ref{homology-lemma-double-complex-gives-resolution}.
\end{proof}

\begin{lemma}
\label{lemma-localize-injective}
Let $\mathcal{C}$ be a site. For $K$ in $\text{SR}(\mathcal{C})$
the functor $j^{-1}$ sends injective abelian sheaves to injective
abelian sheaves. Similarly, the functor $j^{-1}$ sends K-injective
complexes of abelian sheaves to K-injective complexes of
abelian sheaves.
\end{lemma}

\begin{proof}
The first statement is the natural generalization of
Cohomology on Sites, Lemma
\ref{sites-cohomology-lemma-cohomology-of-open}
to semi-representable objects.
In fact, it follows from this lemma
by the product decomposition of $\Sh(\mathcal{C}/K)$
and the description of the functor $j^{-1}$ given above.
The second statement is the natural generalization of
Cohomology on Sites, Lemma
\ref{sites-cohomology-lemma-restrict-K-injective-to-open}
and follows from it by the product decomposition of the topos.

\medskip\noindent
Alternative: since $j$ induces a localization of topoi by
Lemma \ref{lemma-localize-compare} part (1)
it also follows immediately from
Cohomology on Sites, Lemmas \ref{sites-cohomology-lemma-cohomology-of-open}
and \ref{sites-cohomology-lemma-restrict-K-injective-to-open}
by enlarging the site; compare with the proof of
Cohomology on Sites, Lemma
\ref{sites-cohomology-lemma-cohomology-on-sheaf-sets}
in the case of injective sheaves.
\end{proof}

\begin{lemma}
\label{lemma-hypercovering-descent-sheaves}
Let $\mathcal{C}$ be a site with equalizers and fibre products.
Let $K$ be a hypercovering. Then
\begin{enumerate}
\item $a^{-1} : \Sh(\mathcal{C}) \to \Sh((\mathcal{C}/K)_{total})$
is fully faithful with essential image the cartesian sheaves of sets,
\item $a^{-1} : \textit{Ab}(\mathcal{C}) \to
\textit{Ab}((\mathcal{C}/K)_{total})$
is fully faithful with essential image the cartesian sheaves
of abelian groups.
\end{enumerate}
In both cases $a_*$ provides the quasi-inverse functor.
\end{lemma}

\begin{proof}
The case of abelian sheaves follows immediately from the case
of sheaves of sets as the functor $a^{-1}$ commutes with products.
In the rest of the proof we work with sheaves of sets.
Observe that $a^{-1}\mathcal{F}$ is cartesian for
$\mathcal{F}$ in $\Sh(\mathcal{C})$ by
Lemma \ref{lemma-augmentation-cartesian-module}.
It suffices to show that the adjunction map
$\mathcal{F} \to a_*a^{-1}\mathcal{F}$
is an isomorphism $\mathcal{F}$ in $\Sh(\mathcal{C})$
and that for a cartesian sheaf
$\mathcal{G}$ on $(\mathcal{C}/K)_{total}$
the adjunction map
$a^{-1}a_*\mathcal{G} \to \mathcal{G}$ is an isomorphism.

\medskip\noindent
Let $\mathcal{F}$ be a sheaf on $\mathcal{C}$.
Recall that $a_*a^{-1}\mathcal{F}$ is the equalizer
of the two maps $a_{0, *}a_0^{-1}\mathcal{F} \to a_{1, *}a_1^{-1}\mathcal{F}$,
see Lemma \ref{lemma-comparison}.
By Lemma \ref{lemma-push-pull-localization}
$$
a_{0, *}a_0^{-1}\mathcal{F} = \SheafHom(F(K_0)^\#, \mathcal{F})
\quad\text{and}\quad
a_{1, *}a_1^{-1}\mathcal{F} = \SheafHom(F(K_1)^\#, \mathcal{F})
$$
On the other hand, we know that
$$
\xymatrix{
F(K_1)^\# \ar@<1ex>[r] \ar@<-1ex>[r] &
F(K_0)^\# \ar[r] & \text{final object }*\text{ of }\Sh(\mathcal{C})
}
$$
is a coequalizer diagram in sheaves of sets by definition of
a hypercovering. Thus it suffices to prove
that $\SheafHom(-, \mathcal{F})$ transforms coequalizers
into equalizers which is immediate from the construction
in Sites, Section \ref{sites-section-glueing-sheaves}.

\medskip\noindent
Let $\mathcal{G}$ be a cartesian sheaf on $(\mathcal{C}/K)_{total}$.
We will show that $\mathcal{G} = a^{-1}\mathcal{F}$ for some sheaf
$\mathcal{F}$ on $\mathcal{C}$. This will finish the proof because
then $a^{-1}a_*\mathcal{G} = a^{-1}a_*a^{-1}\mathcal{F} =
a^{-1}\mathcal{F} = \mathcal{G}$ by the result of the previous paragraph.
Set $\mathcal{K}_n = F(K_n)^\#$ for $n \geq 0$. Then we have maps of sheaves
$$
\xymatrix{
\mathcal{K}_2
\ar@<1ex>[r]
\ar@<0ex>[r]
\ar@<-1ex>[r]
&
\mathcal{K}_1
\ar@<0.5ex>[r]
\ar@<-0.5ex>[r]
&
\mathcal{K}_0
}
$$
coming from the fact that $K$ is a simplicial semi-representable object.
The fact that $K$ is a hypercovering means that
$$
\mathcal{K}_1 \to \mathcal{K}_0 \times \mathcal{K}_0
\quad\text{and}\quad
\mathcal{K}_2 \to
\left(\text{cosk}_1(
\xymatrix{
\mathcal{K}_1
\ar@<0.5ex>[r]
\ar@<-0.5ex>[r]
&
\mathcal{K}_0 \ar[l]
})\right)_2
$$
are surjective maps of sheaves. Using the description of cartesian sheaves on
$(\mathcal{C}/K)_{total}$ given in Lemma \ref{lemma-characterize-cartesian}
and using the description of $\Sh(\mathcal{C}/K_n)$ in
Lemma \ref{lemma-localize-compare}
we find that our problem can be entirely formulated\footnote{Even though it
does not matter what the precise formulation is, we spell it out:
the problem is to show that given an object
$\mathcal{G}_0/\mathcal{K}_0$ of $\Sh(\mathcal{C})/\mathcal{K}_0$
and an isomorphism
$$
\alpha :
\mathcal{G}_0 \times_{\mathcal{K}_0, \mathcal{K}(\delta^1_1)} \mathcal{K}_1 \to
\mathcal{G}_0 \times_{\mathcal{K}_0, \mathcal{K}(\delta^1_0)} \mathcal{K}_1
$$
over $\mathcal{K}_1$ satisfying a cocycle condition in
$\Sh(\mathcal{C})/\mathcal{K}_2$, there exists
$\mathcal{F}$ in $\Sh(\mathcal{C})$ and an isomorphism
$\mathcal{F} \times \mathcal{K}_0 \to \mathcal{G}_0$ over $\mathcal{K}_0$
compatible with $\alpha$.} in terms of
\begin{enumerate}
\item the topos $\Sh(\mathcal{C})$, and
\item the simplicial object $\mathcal{K}$ in $\Sh(\mathcal{C})$
whose terms are $\mathcal{K}_n$.
\end{enumerate}
Thus, after replacing $\mathcal{C}$ by a different site $\mathcal{C}'$
as in Sites, Lemma \ref{sites-lemma-topos-good-site}, we may assume
$\mathcal{C}$ has all finite limits,
the topology on $\mathcal{C}$ is subcanonical,
a family $\{V_j \to V\}$ of morphisms of $\mathcal{C}$
is a covering if and only if $\coprod h_{V_j} \to V$ is surjective, and
there exists a simplicial object $U$ of $\mathcal{C}$
such that $\mathcal{K}_n = h_{U_n}$ as simplicial sheaves.
Working backwards through the equivalences we may assume
$K_n = \{U_n\}$ for all $n$.

\medskip\noindent
Let $X$ be the final object of $\mathcal{C}$.
Then $\{U_0 \to X\}$ is a covering,
$\{U_1 \to U_0 \times U_0\}$ is a covering, and
$\{U_2 \to (\text{cosk}_1 \text{sk}_1 U)_2\}$ is a covering.
Let us use $d^n_i : U_n \to U_{n - 1}$ and
$s^n_j : U_n \to U_{n + 1}$ the morphisms corresponding
to $\delta^n_i$ and $\sigma^n_j$ as in
Simplicial, Definition \ref{simplicial-definition-face-degeneracy}.
By abuse of notation, given a morphism
$c : V \to W$ of $\mathcal{C}$ we denote the morphism of topoi
$c : \Sh(\mathcal{C}/V) \to \Sh(\mathcal{C}/W)$ by the same letter.
Now $\mathcal{G}$ is given by a sheaf $\mathcal{G}_0$
on $\mathcal{C}/U_0$ and an isomorphism
$\alpha : (d^1_1)^{-1}\mathcal{G}_0 \to (d^1_0)^{-1}\mathcal{G}_0$
satisfying the cocycle condition on $\mathcal{C}/U_2$
formulated in Lemma \ref{lemma-characterize-cartesian}.
Since $\{U_2 \to (\text{cosk}_1 \text{sk}_1 U)_2\}$
is a covering, the corresponding pullback functor
on sheaves is faithful (small detail omitted).
Hence we may replace $U$ by $\text{cosk}_1 \text{sk}_1 U$, because
this replaces $U_2$ by $(\text{cosk}_1 \text{sk}_1 U)_2$ and leaves
$U_1$ and $U_0$ unchanged. Then
$$
(d^2_0, d^2_1, d^2_2) : U_2 \to U_1 \times U_1 \times U_1
$$
is a monomorphism whose its image on $T$-valued points is
described in Simplicial, Lemma \ref{simplicial-lemma-work-out}.
In particular, there is a morphism $c$ fitting into a commutative diagram
$$
\xymatrix{
U_1 \times_{(d^1_1, d^1_0), U_0 \times U_0, (d^1_1, d^1_0)} U_1
\ar[d] \ar[rr]_c & & U_2 \ar[d] \\
U_1 \times U_1
\ar[rr]^{(\text{pr}_1, \text{pr}_2, s^0_0 \circ d^1_1 \circ \text{pr}_1)} & &
U_1 \times U_1 \times U_1
}
$$
as going around the other way defines a point of $U_2$.
Pulling back the cocycle condition for $\alpha$ on $U_2$
translates into the condition that the pullbacks of $\alpha$
via the projections to
$U_1 \times_{(d^1_1, d^1_0), U_0 \times U_0, (d^1_1, d^1_0)} U_1$
are the same as the pullback of $\alpha$ via
$s^0_0 \circ d^1_1 \circ \text{pr}_1$ is the identity map
(namely, the pullback of $\alpha$ by $s^0_0$ is the identity).
By Sites, Lemma \ref{sites-lemma-glue-maps}
this means that $\alpha$ comes from an isomorphism
$$
\alpha' : \text{pr}_1^{-1}\mathcal{G}_0 \to \text{pr}_2^{-1}\mathcal{G}_0
$$
of sheaves on $\mathcal{C}/U_0 \times U_0$.
Then finally, the morphism $U_2 \to U_0 \times U_0 \times U_0$
is surjective on associated sheaves as is easily seen using the
surjectivity of $U_1 \to U_0 \times U_0$
and the description of $U_2$ given above. Therefore $\alpha'$
satisfies the cocycle condition on $U_0 \times U_0 \times U_0$.
The proof is finished by an application of
Sites, Lemma \ref{sites-lemma-mapping-property-glue}
to the covering $\{U_0 \to X\}$.
\end{proof}

\begin{lemma}
\label{lemma-hypercovering-cech-complex}
Let $\mathcal{C}$ be a site with equalizers and fibre products.
Let $K$ be a hypercovering. The {\v C}ech complex
of Lemma \ref{lemma-augmentation-cech-complex} associated to
$a^{-1}\mathcal{F}$
$$
a_{0, *}a_0^{-1}\mathcal{F} \to a_{1, *}a_1^{-1}\mathcal{F} \to
a_{2, *}a_2^{-1}\mathcal{F} \to \ldots
$$
is equal to the complex $\SheafHom(s(\mathbf{Z}_{F(K)}^\#), \mathcal{F})$.
Here $s(\mathbf{Z}_{F(K)}^\#)$ is as in
Hypercoverings, Definition \ref{hypercovering-definition-homology}.
\end{lemma}

\begin{proof}
By Lemma \ref{lemma-push-pull-localization} we have
$$
a_{n, *}a_n^{-1}\mathcal{F} = \SheafHom'(F(K_n)^\#, \mathcal{F})
$$
where $\SheafHom'$ is as in Sites, Section \ref{sites-section-glueing-sheaves}.
The boundary maps in the complex of
Lemma \ref{lemma-augmentation-cech-complex}
come from the simplicial structure.
Thus the equality of complexes comes 
from the canonical identifications
$\SheafHom'(\mathcal{G}, \mathcal{F}) =
\SheafHom(\mathbf{Z}_\mathcal{G}, \mathcal{F})$ for
$\mathcal{G}$ in $\Sh(\mathcal{C})$.
\end{proof}

\begin{proof}
First, let $\mathcal{I}$ be an injective abelian sheaf on $\mathcal{C}$.
Then the spectral sequence of
Lemma \ref{lemma-augmentation-spectral-sequence}
for the sheaf $a^{-1}\mathcal{I}$ degenerates as
$(a^{-1}\mathcal{I})_p = a_p^{-1}\mathcal{I}$
is injective by Lemma \ref{lemma-localize-injective}.
Thus the complex
$$
a_{0, *}a_0^{-1}\mathcal{I} \to
a_{1, *}a_1^{-1}\mathcal{I} \to
a_{2, *}a_2^{-1}\mathcal{I} \to\ldots
$$
computes $Ra_*a^{-1}\mathcal{I}$. By
Lemma \ref{lemma-hypercovering-cech-complex}
this is equal to the complex
$\SheafHom(s(\mathbf{Z}_{F(K)}^\#), \mathcal{I})$.
Because $K$ is a hypercovering, we see that
$s(\mathbf{Z}_{F(K)}^\#)$ is exact in degrees $> 0$ by
Hypercoverings, Lemma \ref{hypercovering-lemma-acyclic-hypercover-sheaves}
applied to the simplicial presheaf $F(K)$.
Since $\mathcal{I}$ is injective, the functor $\SheafHom(-, \mathcal{I})$
is exact and we conclude that
$\SheafHom(s(\mathbf{Z}_{F(K)}^\#), \mathcal{I})$
is exact in positive degrees. We conclude that
$R^pa_*a^{-1}\mathcal{I} = 0$ for $p > 0$.
On the other hand, we have $\mathcal{I} = a_*a^{-1}\mathcal{I}$
by Lemma \ref{lemma-hypercovering-descent-sheaves}.

\medskip\noindent
Bounded case. Let $E \in D^+(\mathcal{C})$.
Choose a bounded below complex $\mathcal{I}^\bullet$ of injectives
representing $E$. By the result of the first paragraph and
Leray's acyclicity lemma
(Derived Categories, Lemma \ref{derived-lemma-leray-acyclicity})
$Ra_*a^{-1}\mathcal{I}^\bullet$
is computed by the complex
$a_*a^{-1}\mathcal{I}^\bullet = \mathcal{I}^\bullet$
and we conclude the lemma is true in this case.

\medskip\noindent
Unbounded case. We urge the reader to skip this, since the argument
is the same as above, except that we use explicit representation
by double complexes to get around convergence issues.
Let $E \in D(\mathcal{C})$.
To show the map $E \to Ra_*a^{-1}E$ is an isomorphism,
it suffices to show for every object $U$ of $\mathcal{C}$ that
$$
R\Gamma(U, E) = R\Gamma(U, Ra_*a^{-1}E)
$$
We will compute both sides and show the map $E \to Ra_*a^{-1}E$
induces an isomorphism. Choose a K-injective
complex $\mathcal{I}^\bullet$ representing $E$. Choose a quasi-isomorphism
$a^{-1}\mathcal{I}^\bullet \to \mathcal{J}^\bullet$
for some K-injective complex $\mathcal{J}^\bullet$ on
$(\mathcal{C}/K)_{total}$. We have
$$
R\Gamma(U, E) = R\Hom(\mathbf{Z}_U^\#, E)
$$
and
$$
R\Gamma(U, Ra_*a^{-1}E) = R\Hom(\mathbf{Z}_U^\#, Ra_*a^{-1}E) =
R\Hom(a^{-1}\mathbf{Z}_U^\#, a^{-1}E)
$$
By Lemma \ref{lemma-simplicial-resolution-augmentation}
we have a quasi-isomorphism
$$
\Big(\ldots \to
g_{2!}(a_2^{-1}\mathbf{Z}_U^\#) \to
g_{1!}(a_1^{-1}\mathbf{Z}_U^\#) \to
g_{0!}(a_0^{-1}\mathbf{Z}_U^\#)\Big)
\longrightarrow
a^{-1}\mathbf{Z}_U^\#
$$
Hence $R\Hom(a^{-1}\mathbf{Z}_U^\#, a^{-1}E)$ is equal to
$$
R\Gamma((\mathcal{C}/K)_{total},
R\SheafHom(
\ldots \to
g_{2!}(a_2^{-1}\mathbf{Z}_U^\#) \to
g_{1!}(a_1^{-1}\mathbf{Z}_U^\#) \to
g_{0!}(a_0^{-1}\mathbf{Z}_U^\#),
\mathcal{J}^\bullet))
$$
By the construction in Cohomology on Sites, Section
\ref{sites-cohomology-section-internal-hom}
and since $\mathcal{J}^\bullet$ is K-injective, we see that
this is represented by the complex of abelian groups with terms
$$
\prod\nolimits_{p + q = n}
\Hom(g_{p!}(a_p^{-1}\mathbf{Z}_U^\#), \mathcal{J}^q) =
\prod\nolimits_{p + q = n}
\Hom(a_p^{-1}\mathbf{Z}_U^\#, g_p^{-1}\mathcal{J}^q)
$$
See Cohomology on Sites, Lemmas
\ref{sites-cohomology-lemma-RHom-into-K-injective} and
\ref{sites-cohomology-lemma-section-RHom-over-U} for more information.
Thus we find that $R\Gamma(U, Ra_*a^{-1}E)$ is computed by
the product total complex $\text{Tot}_\pi(B^{\bullet, \bullet})$
with $B^{p, q} = \Hom(a_p^{-1}\mathbf{Z}_U^\#, g_p^{-1}\mathcal{J}^q)$.
For the other side we argue similarly. First we note that
$$
s(\mathbf{Z}_{F(K)}^\#) \longrightarrow \mathbf{Z}
$$
is a quasi-isomorphism of complexes on $\mathcal{C}$
by Hypercoverings, Lemma \ref{hypercovering-lemma-acyclic-hypercover-sheaves}.
Since $\mathbf{Z}_U^\#$ is a flat sheaf of $\mathbf{Z}$-modules
we see that
$$
s(\mathbf{Z}_{F(K)}^\#) \otimes_\mathbf{Z} \mathbf{Z}_U^\#
\longrightarrow
\mathbf{Z}_U^\#
$$
is a quasi-isomorphism. Therefore
$R\Hom(\mathbf{Z}_U^\#, E)$ is equal to
$$
R\Gamma(\mathcal{C}, R\SheafHom(
s(\mathbf{Z}_{F(K)}^\#) \otimes_\mathbf{Z} \mathbf{Z}_U^\#,
\mathcal{I}^\bullet))
$$
By the construction of $R\SheafHom$ and since $\mathcal{I}^\bullet$
is K-injective, this is represented by the complex of abelian groups
with terms
$$
\prod\nolimits_{p + q = n}
\Hom(\mathbf{Z}^\#_{K_p} \otimes_\mathbf{Z} \mathbf{Z}_U^\#, \mathcal{I}^q)
=
\prod\nolimits_{p + q = n}
\Hom(a_p^{-1}\mathbf{Z}_U^\#, a_p^{-1}\mathcal{I}^q)
$$
The equality of terms follows from the fact that
$\mathbf{Z}^\#_{K_p} \otimes_\mathbf{Z} \mathbf{Z}_U^\# =
a_{p!}a_p^{-1}\mathbf{Z}_U^\#$ by Modules on Sites,
Remark \ref{sites-modules-remark-j-shriek-tensor}.
Thus we find that $R\Gamma(U, E)$ is computed by
the product total complex $\text{Tot}_\pi(A^{\bullet, \bullet})$
with $A^{p, q} = \Hom(a_p^{-1}\mathbf{Z}_U^\#, a_p^{-1}\mathcal{I}^q)$.

\medskip\noindent
Since $\mathcal{I}^\bullet$ is K-injective we see that
$a_p^{-1}\mathcal{I}^\bullet$ is K-injective, see
Lemma \ref{lemma-localize-injective}.
Since $\mathcal{J}^\bullet$ is K-injective we see that
$g_p^{-1}\mathcal{J}^\bullet$ is K-injective, see
Lemma \ref{lemma-restriction-injective-to-component-site}.
Both represent the object $a_p^{-1}E$.
Hence for every $p \geq 0$ the map of complexes
$$
A^{p, \bullet} = \Hom(a_p^{-1}\mathbf{Z}_U^\#, a_p^{-1}\mathcal{I}^\bullet)
\longrightarrow
\Hom(a_p^{-1}\mathbf{Z}_U^\#, g_p^{-1}\mathcal{J}^\bullet) = B^{p, \bullet}
$$
induced by $g_p^{-1}$ applied to the given map
$a^{-1}\mathcal{I}^\bullet \to \mathcal{J}^\bullet$
is a quasi-isomorphisms as these complexes both compute
$$
R\Hom(a_p^{-1}\mathbf{Z}_U^\#, a_p^{-1}E)
$$
By More on Algebra, Lemma \ref{more-algebra-lemma-prod-qis-gives-qis}
we conclude that the right vertical arrow in the commutative diagram
$$
\xymatrix{
R\Gamma(U, E) \ar[r] \ar[d] &
\text{Tot}_\pi(A^{\bullet, \bullet}) \ar[d] \\
R\Gamma(U, Ra_*a^{-1}E) \ar[r] &
\text{Tot}_\pi(B^{\bullet, \bullet})
}
$$
is a quasi-isomorphism. Since we saw above that the horizontal arrows
are quasi-isomorphisms, so is the left vertical arrow.
\end{proof}

\begin{lemma}
\label{lemma-acyclic-hypercover-sheaves}
Let $\mathcal{C}$ be a site.
Let $K$ be a simplicial presheaf.
Let $\mathcal{G}$ be a presheaf.
Let $K \to \mathcal{G}$ be an augmentation of $K$
towards $\mathcal{G}$. Assume that
\begin{enumerate}
\item The morphism of presheaves $K_0 \to \mathcal{G}$ becomes
a surjection after sheafification.
\item The morphism
$$
(d^1_0, d^1_1) :
K_1
\longrightarrow
K_0 \times_\mathcal{G} K_0
$$
becomes a surjection after sheafification.
\item For every $n \geq 1$ the morphism
$$
K_{n + 1} \longrightarrow (\text{cosk}_n \text{sk}_n K)_{n + 1}
$$
turns into a surjection after sheafification.
\end{enumerate}
Then $H_i(K) = 0$ for $i > 0$ and
$H_0(K) = \mathbf{Z}_\mathcal{G}^\#$.
\end{lemma}

\begin{proof}
Denote $K^n = \text{cosk}_n \text{sk}_n K$ for $n \geq 1$.
Define $K^0$ as the simplicial object with terms
$(K^0)_n$ equal to the $(n + 1)$-fold fibred product
$K_0 \times_\mathcal{G} \ldots \times_\mathcal{G} K_0$,
see Simplicial,
Example \ref{simplicial-example-fibre-products-simplicial-object}.
We have morphisms
$$
K \longrightarrow \ldots \to K^n \to K^{n - 1} \to \ldots \to K^1 \to K^0.
$$
The morphisms $K \to K^i$, $K^j \to K^i$ for $j \geq i \geq 1$ come
from the universal properties of the $\text{cosk}_n$ functors.
The morphism $K^1 \to K^0$ is the canonical morphism
from
Simplicial, Remark \ref{simplicial-remark-augmentation}.
We also recall that $K^0 \to \text{cosk}_1 \text{sk}_1 K^0$
is an isomorphism, see
Simplicial, Lemma \ref{simplicial-lemma-cosk-minus-one}.

\medskip\noindent
By Lemma \ref{lemma-compare-cosk0} we see that
$H_i(K^0) = 0$ for $i > 0$ and $H_0(K^0) = \mathbf{Z}_\mathcal{G}^\#$.

\medskip\noindent
Pick $n \geq 1$. Consider the morphism $K^n \to K^{n - 1}$.
It is an isomorphism on terms of degree $< n$.
Note that $K^n \to \text{cosk}_n \text{sk}_n K^n$ and
$K^{n - 1} \to \text{cosk}_n \text{sk}_n K^{n - 1}$
are isomorphisms. Note that $(K^n)_n = K_n$ and
that $(K^{n - 1})_n = (\text{cosk}_{n - 1} \text{sk}_{n - 1} K)_n$.
Hence by assumption, we have that $(K^n)_n \to (K^{n - 1})_n$
is a morphism of presheaves which becomes surjective after
sheafification. By Lemma \ref{lemma-acyclicity} we conclude that
$H_i(K^n) = H_i(K^{n - 1})$.
Combined with the above this proves the lemma.
\end{proof}

\bibliographystyle{plain}
\bibliography{my}
\end{document}
